
%----------------------------------------------------------------------------------------
%	                            PORTADA
%----------------------------------------------------------------------------------------
% \makeportada{}{Daniel Carbajales Soto}{ 1° CFGS Automatización y Robótica Industrial}{}

%----------------------------------------------------------------------------------------
%                               TABLE OF CONTENTS
%----------------------------------------------------------------------------------------
% \maketoc
%----------------------------------------------------------------------------------------
%                                    HEADER 
%----------------------------------------------------------------------------------------
% \header{}{Daniel Carbajales Soto}{1° CFGSARI}
%----------------------------------------------------------------------------------------
%	                            TABLES
%----------------------------------------------------------------------------------------
% this is just a demo in case i need to make a table
%
%
%
% \begin{tabularx}{\linewidth}{|C{2cm}|Y|Y|C{2cm}|}
% \hline
% \makecell{Header 1} & \makecell{Header 2 \\ with linebreak} & \makecell{Header 3} & \makecell{Header 4} \\
% \hline
% Row 1 & This is a long cell text that automatically wraps in the Y column & Another long cell & Short \\
% \hline
% Row 2 & \multirow{2}{*}{Multirow text} & More text & Short \\
% \cline{1-1} \cline{3-4}
% Row 3 & & Extra text & Short \\
% \hline
% \end{tabularx}


%----------------------------------------------------------------------------------------
%	                            DOCUMENT
%----------------------------------------------------------------------------------------
\phantomsection
\addcontentsline{toc}{chapter}{Ej 1}
\fsection{\boldit{\textcolor{waveBlue2}{INVESTGA}}}
Haz una búsqueda en internet y localiza la página web de UNILEVER, tanto global como nacional, 
y extrae algunas conclusiones sobre su apuesta por la sostenibilidad y la mejora del planeta.


\fsubsection{Compromisos Globales de Unilever}
Unilever ha establecido un Plan de Acción de Crecimiento Sostenible con 15 metas ambiciosas a corto y mediano plazo,
enfocadas en cuatro áreas prioritarias: \boldit{clima, naturaleza, plásticos y medios de vida.} \\
\begin{itemize}
	\item Clima: prioriza acciones
		urgentes para reducir emisiones y adaptarse a desafíos globales.
	\item Naturaleza: trabaja en la conservación y restauración 
		de ecosistemas. 
	\item Plásticos: busca reducir el uso de plástico virgen (logrando un 23\% de disminución desde 2019 e
		incorporando un 21\% de plástico reciclado en envases) y mejorar el reciclaje. 
	\item Medios de vida: 
		apoya comunidades vulnerables mediante prácticas regenerativas que cubren el 50\% de su huella 
		territorial en alimentos. La compañía se centra en la influencia sistémica, colaborando con gobiernos y 
		sociedades diversas para eliminando barreras e integrando la sostenibilidad en toda la organización; 
		logrando avances iniciales alentadores.
\end{itemize}

\fsubsection{Iniciativas en España}
En España, Unilever alinea sus esfuerzos con los globales, avanzando en objetivos climáticos:
\begin{itemize}
	\item Cero emisiones en sus fábricas para 2030.
	\item Protección y restauración de un millón de ecosistemas naturales globales.
\end{itemize}
Sus principales iniciativas incluyen:
\begin{itemize}
	\item La transición a envases sostenibles 
	\item Prácticas regenerativas en agricultura.
\end{itemize} 
Con estas iniciativas que contribuyen a la Agenda 2030. Observamos un enfoque pragmático reduciendo el impacto local, como en 
el sector alimentario, promoviendo empleo verde y eficiencia energética.

\fsubsection{Conclusiones}
Unilever demuestra una apuesta estratégica y holística por la sostenibilidad, pasando de medidas tácticas a transformaciones sistémicas que generan valor compartido. Su enfoque no solo mitiga riesgos ambientales, sino que impulsa innovación y resiliencia, posicionándola como líder en la mejora del planeta. Sin embargo, el éxito depende de colaboraciones externas para escalar impactos.






%%%%%%%%%%%%%%%%%%%%%%%%%%%%%%%%%%%%%%%%%%%%%%%%%%%%%%%%%%
%%%%%%%%%%%%%%%%%%%% bibliography %%%%%%%%%%%%%%%%%%%%%%%%
% \nocite{*}                                 %%%%%%%%%%%%%
% \printbibliography                         %%%%%%%%%%%%%
%%%%%%%%%%%%%%%%%%%%%%%%%%%%%%%%%%%%%%%%%%%%%%%%%%%%%%%%%%
%%%%%% Last page in case there is no bibliography %%%%%%%%
% \label{Last Page}                           %%%%%%%%%%%%%
%%%%%%%%%%%%%%%%%%%%%%%%%%%%%%%%%%%%%%%%%%%%%%%%%%%%%%%%%%

